\section{Preliminaries} \label{sec:prelim}


\subsection[The first-order jet bundle $J^{1}\pi$]{The f\/irst-order jet bundle $\boldsymbol{J^{1}\pi}$} \label{subsec:fn}

The appropriate geometric setting for the study of time-dependent
SODE is the af\/f\/ine jet bundle $(J^{1}\pi,\pi_{10},M)$~\cite{saunders89}.
We consider an $(n+1)$-dimensional, real, smooth manifold $M$, which
is f\/ibred over $\mathbb{R}$, $\pi:M\rightarrow\mathbb{R}$, and represents
the space-time. The f\/irst jet bundle of $\pi$ is denoted by $\pi_{10}:J^{1}\pi\to M$,
$\pi_{10}(j_{t}^{1}\phi)=\phi(t),$ for $\phi$ a local section of
$\pi$ and $j_{t}^{1}\phi$ the f\/irst-order jet of~$\phi$ at~$t$.
A local coordinate system $(t,x^{i})_{i\in \{1,\dots,n\}}$ on $M$
induces a local coordinate system on $J^{1}\pi$, denoted by $(t,x^{i},y^{i})$.
Submersion $\pi_{10}$ induces a \emph{natural foliation} on $J^{1}\pi$
such that~$(t,x^{i})$ are transverse coordinates for this foliation,
while~$(y^{i})$ are coordinates for the leaves of the foliation.
Throughout the paper we consider Latin indices $i\in \{1,\dots,n\}$
and Greek indices $\alpha\in \{0,\dots,n\}$, using the notation~$\left(x^{\alpha}\right)=(t=x^{0},x^{i})$.

In this article we use the Fr\"olicher--Nijenhuis theory \cite{frolicher56,grifone00,KMS93}
of derivations of vector-valued dif\/ferential forms on the f\/irst jet
bundle $J^{1}\pi$. We adopt the following notations: $C^{\infty}(J^{1}\pi)$
for the ring of smooth functions on~$J^{1}\pi$, $\mathfrak{X}(J^{1}\pi)$
for the~$C^{\infty}$ module of vector f\/ields on $J^{1}\pi$ and $\Lambda^{k}(J^{1}\pi)$
for the~$C^{\infty}$ module of $k$-forms on~$J^{1}\pi$. The $C^{\infty}$
module of $(r,s)$-type vector f\/ields on~$J^{1}\pi$ is denoted by
$\mathcal{T}_{s}^{r}(J^{1}\pi)$ and the tensor algebra on~$J^{1}\pi$
is denoted by $\mathcal{T}(J^{1}\pi)$. The graded algebra of dif\/ferential
forms on $J^{1}\pi$ is written as $\Lambda(J^{1}\pi)=\bigoplus_{k\in \{1,\dots,2n+1\}}\Lambda^{k}(J^{1}\pi)$.
We denote by~$S^{k}(J^{1}\pi)$ the space of symmetric $(0,k)$ tensors
on~$J^{1}\pi$ and by $\Psi(J^{1}\pi)=\bigoplus_{k\in \{1,\dots,2n+1\}}\Psi^{k}(J^{1}\pi)$
the graded algebra of vector-valued dif\/ferential forms on $J^{1}\pi$.
Throughout the paper we assume that all objects are $C^{\infty}$-smooth
where def\/ined.

A parametrized curve on $M$ is a section of $\pi$: $\gamma:\mathbb{R}\rightarrow M$,
$\gamma(t)=(t,x^{i}(t))$. Its \emph{first-order jet prolongation}
$J^{1}\gamma:t\in\mathbb{R}\rightarrow J^{1}\gamma(t)=\left(t,x^{i}(t),dx^{i}/dt\right)\in J^{1}\pi$
is a section of the f\/ibration $\pi_{1}:=\pi\circ\pi_{10}:  J^{1}\pi\rightarrow\mathbb{R}$.

Let $VJ^{1}\pi$ be the \emph{vertical subbundle} of $TJ^{1}\pi,$
$VJ^{1}\pi=\{\xi\in TJ^{1}\pi,\,  D\pi_{10}(\xi)=0\}\subset TJ^{1}\pi$.
The f\/ibers $V_{u}J^{1}\pi=\operatorname{Ker} D_{u}\pi_{10}$, $u\in J^{1}\pi$ determine
a regular, $n$-dimensional, integrable vertical distribution. Remark
that $V_{u}J^{1}\pi=\operatorname{spann}\{\partial/\partial y^{i}\}$
and its annihilators are the \emph{contact $1$-forms} $\delta x^{i}=dx^{i}-y^{i}dt$, $i\in \{1,\dots,n\}$
and \emph{basic $1$-forms} $\lambda dt$, $\lambda\in C^{\infty}(J^{1}\pi)$.
The \emph{vertical endomorphism} $J=\frac{\partial}{\partial y^{i}}\otimes\delta x^{i}$
is a vector-valued 1-form on~$J^{1}\pi$, with $\operatorname{Im}J=V(J^{1}\pi)$,
$V(J^{1}\pi)\subset\operatorname{Ker}J$ and $J^{2}=0$.

Its Fr\"olicher--Nijenhuis tensor is given by
\begin{gather*}
N_{J}=\frac{1}{2}[J,J]=-\frac{\partial}{\partial y^{i}}\otimes\delta x^{i}\wedge dt=-J\wedge dt.
%\label{eq:NijenhuisJ}
\end{gather*}
 Consequently, $d_{J}^{2}=d_{N_{J}}=-d_{J\wedge dt}\neq0$ and therefore
$d_{J}$-exact forms on $J^{1}\pi$ may not be $d_{J}$-closed. Here
$d_{J}$ is the \emph{exterior derivative} with respect to the vertical
endomorphisms.

\begin{remark}
For $A\in\Psi^{1}(J^{1}\pi)$ a vector-valued 1-form, the exterior
derivative with respect to $A$ is a derivation of degree 1 given
by $d_{A}=i_{A}\circ d-d\circ i_{A}.$
\end{remark}
A $k$-form $\omega$ on $J^{1}\pi$, $k\geq1$, is called \emph{semi-basic}
if it vanishes whenever one of the arguments is vertical.

A vector-valued $k$-form $A$ on $J^{1}\pi$ is called semi-basic
if it takes values in the vertical bundle and it vanishes whenever
one of the arguments is vertical.

A semi-basic $k$-form satisf\/ies the relation $i_{J}\theta=0$ and
locally can be expressed as $\theta=\theta_{0}dt+\theta_{i}\delta x^{i}$.
For example, contact 1-forms $\delta x^{i}$ are semi-basic 1-forms.

If a vector-valued $k$-form $A$ is semi-basic, then $J\circ A=0$
and $i_{J}A=0$. The vertical endomorphism $J$ is a vector-valued,
semi-basic $1$ -form.

Locally, a semi-basic $k$-form $\theta$ has the next form \begin{eqnarray*}
\theta=\frac{1}{k!}\theta_{i_{1}\dots i_{k}}(x^{\alpha},y^{j})\delta x^{i_{1}}\wedge\cdots\wedge\delta x^{i_{k}}+\frac{1}{(k-1)!}\widetilde{\theta}_{i_{1}\dots i_{k-1}}(x^{\alpha},y^{j})\delta x^{i_{1}}\wedge\cdots\wedge\delta x^{i_{k-1}}\wedge dt.\end{eqnarray*}


For simplicity, we denote by $T^{*}$ the vector bundle of $1$-forms
on $J^{1}\pi$, by $T_{v}^{*}$ the vector bundle of semi-basic $1$-forms
on $J^{1}\pi$ and by $\Lambda^{k}T_{v}^{*}$ the vector bundle of
semi-basic $k$-forms on $J^{1}\pi$. We also denote by $\Lambda_{v}^{k}=\operatorname{Sec}\left(\Lambda^{k}T_{v}^{*}\right)$
the $C^{\infty}(J^{1}\pi)$-module of sections of $\Lambda^{k}T_{v}^{*}$
and by $S^{k}T^{*}$ the vector bundle of symmetric tensors of $(0,k)$-type
on $J^{1}\pi$. $S^{1}T^{*}$ will be identif\/ied with $T^{*}$.

A \emph{semispray} is a globally def\/ined vector f\/ield $S$ on $J^{1}\pi$
such that \begin{gather*}
J(S)=0\qquad\textrm{and}\qquad dt(S)=1.%\label{eq:4}
\end{gather*}
 The integral curves of a semispray are f\/irst-order jet prolongations
of sections of $\pi\circ\pi_{10}:\!  J^{1}\pi\!\rightarrow\!\mathbb{R}$.
Locally, a semispray has the form{\samepage \begin{gather}
S=\frac{\partial}{\partial t}+y^{i}\frac{\partial}{\partial x^{i}}-2G^{i}(x^{\alpha},y^{j})\frac{\partial}{\partial y^{i}},\label{eq:5}
\end{gather}
 where functions $G^{i}$, called the semispray coef\/f\/icients, are
locally def\/ined on~$J^{1}\pi$.}

A parametrized curve $\gamma:  I\rightarrow M$ is a \emph{geodesic}
of $S$ if $S\circ J^{1}\gamma=\frac{d}{dt}(J^{1}\gamma).$

In local coordinates, $\gamma(t)=(t,x^{i}(t))$ is a geodesic of the
semispray $S$ given by \eqref{eq:5} if and only if it satisf\/ies
the system of SODE \begin{gather}
\frac{d^{2}x^{i}}{dt^{2}}+2G^{i}\left(t,x,\frac{dx}{dt}\right)=0.\label{sode}
\end{gather}


Therefore such a system of time-dependent SODE can be identif\/ied with
a semispray on $J^{1}\pi$.

{\bf Canonical nonlinear connection.}
A~\emph{nonlinear connection} on $J^{1}\pi$ is an $(n+1)$-dimensional
distribution $H:u\in J^{1}\pi\mapsto H_{u}\subset T_{u}J^{1}\pi$,
supplementary to $VJ^{1}\pi$: $\forall\, u\in J^{1}\pi$, $T_{u}J^{1}\pi=H_{u}\oplus V_{u}$.

A semispray $S$ induces a nonlinear connection on $J^{1}\pi$, given
by the \emph{almost product structure} $\Gamma=-\mathcal{L}_{S}J+S\otimes dt$, $\Gamma^{2}={\rm Id}$.
The \emph{horizontal projector} that corresponds to this almost product
structure is $h=\frac{1}{2}\left({\rm Id}-\mathcal{L}_{S}J+S\otimes dt\right)$
and the \emph{vertical projector} is $v={\rm Id}-h.$

The horizontal subspace is spanned by $S$ and by $\frac{\delta}{\delta x^{i}}:=\frac{\partial}{\partial x^{i}}-N_{i}^{j}\frac{\partial}{\partial y^{j}}$, where $N_{j}^{i}=\frac{\partial G^{i}}{\partial y^{j}}$.
In this paper we prefer to work with the following adapted basis and
cobasis:
\begin{gather}
\left\{ S,\frac{\delta}{\delta x^{i}},\frac{\partial}{\partial y^{i}}\right\} ,\qquad\{dt,\delta x^{i},\delta y^{i}\},\label{eq:11}
\end{gather}
 with $\delta x^{i}$ the \emph{contact $1$-forms} and $\delta y^{i}=dy^{i}+N_{\alpha}^{i}dx^{\alpha}$, $N_{0}^{i}=2G^{i}-N_{j}^{i}y^{j}$.
Functions $N_{j}^{i}$ and $N_{0}^{i}$ are the coef\/f\/icients of the
nonlinear connection induced by the semispray~$S$.

With respect to basis and cobasis \eqref{eq:11}, the horizontal and
vertical projectors are locally expressed as $h=S\otimes dt+\frac{\delta}{\delta x^{i}}\otimes\delta x^{i}$, $v=\frac{\partial}{\partial y^{i}}\otimes\delta y^{i}$.
We consider the $(1,1)$-type tensor f\/ield $\mathbb{F}=h\circ\mathcal{L}_{S}h-J$,
which corresponds to the almost complex structure in the autonomous
case. It satisf\/ies $\mathbb{F}^{3}+\mathbb{F}=0,$ which means that
it is an~$f(3,1)$ structure. It can be expressed locally as $\mathbb{F}=\frac{\delta}{\delta x^{i}}\otimes\delta y^{i}-\frac{\partial}{\partial y^{j}}\otimes\delta x^{i}$.



{\bf Curvature.}
The following properties for the torsion and curvature of the nonlinear
connection induced by the semispray are proved in~\cite{B-C-2010}.

The weak torsion tensor f\/ield of the nonlinear connection $\Gamma$
vanishes: $[J,h]=0,$ which is equivalent also with $[J,\Gamma]=0$.

The curvature tensor $R=N_{h}$ of the nonlinear connection $\Gamma$
is a vector-valued semi-basic 2-form, locally given by
\begin{gather}
R=\frac{1}{2}[h,h]=\frac{1}{2}R_{ij}^{k}\frac{\partial}{\partial y^{k}}\otimes\delta x^{i}\wedge\delta x^{j}+R_{i}^{j}\frac{\partial}{\partial y^{j}}\otimes dt\wedge\delta x^{i},\label{eq:19}
\end{gather}
 where \begin{gather*}
R_{jk}^{i}=\frac{\delta N_{j}^{i}}{\delta x^{k}}-\frac{\delta N_{k}^{i}}{\delta x^{j}}%\label{eq:Rijk}
\end{gather*}
 and \begin{gather}
R_{j}^{i}=2\frac{\partial G^{i}}{\partial x^{j}}-\frac{\partial G^{i}}{\partial y^{k}}\frac{\partial G^{k}}{\partial y^{j}}-S\left(\frac{\partial G^{i}}{\partial y^{j}}\right).\label{eq:local-Jacobi}
\end{gather}


The \emph{Jacobi endomorphism} is def\/ined as \begin{gather}
\Phi=v\circ\mathcal{L}_{S}h=\mathcal{L}_{S}h-\mathbb{F}-J.\label{eq:22}\end{gather}
 Jacobi endomorphism $\Phi$ is a semi-basic, vector-valued 1-form
and satisf\/ies $\Phi^{2}=0$. Locally, can be expressed as $\Phi=R_{i}^{j}\frac{\partial}{\partial y^{j}}\otimes\delta x^{i},$
where $R_{j}^{i}$ are given by~(\ref{eq:local-Jacobi}).

The Jacobi endomorphism and the curvature of the nonlinear connection
are related by the following formulae:
\begin{gather}
\Phi   =   i_{S}R,\label{eq:23'}\\
\left[J,\Phi\right]   =   3R+\Phi\wedge dt.\label{eq:24'}
\end{gather}


Remark that $R=0$ if and only if $\Phi=0$.

\begin{definition}
A semispray $S$ is called \emph{isotropic} if its Jacobi endomorphism
has the form \begin{gather}
\Phi=\lambda J,
\label{eq:phi_iso}
\end{gather}
 where $\lambda\in C^{\infty}(J^{1}\pi)$.
\end{definition}
Next we express the isotropy condition \eqref{eq:phi_iso} for a semispray
in terms of the curvature tensor~$R$.

\begin{proposition}
\label{prop:iso} A semispray $S$ is isotropic if and only if its
curvature tensor $R$ has the form
\begin{gather*}
R=\alpha\wedge J,%\label{eq:r_iso}
\end{gather*}
 where $\alpha$ is a semi-basic $1$-form on $J^{1}\pi$.
\end{proposition}

\begin{proof}
Suppose that $S$ is an isotropic SODE. Then there exists $\lambda\in C^{\infty}(J^{1}\pi)$
such that $\Phi=\lambda J$. From (\ref{eq:24'}) it results
\begin{gather*}
3R   =   [J,\lambda J]-\Phi\wedge dt,\\
{}[J,\lambda J]   =   \left(d_{J}\lambda\right)\wedge J-d\lambda\wedge J^{2}+\lambda[J,J] \ \Rightarrow \
R   =   \frac{1}{3}\left(d_{J}\lambda\right)\wedge J-\lambda J\wedge dt=\alpha\wedge J,
\end{gather*}
 with $\alpha=\frac{1}{3}d_{J}\lambda+\lambda dt\in T_{v}^{*}$.

In the above calculus we used the formula \cite{grifone00}\[
[K,gL]= (d_{K}g )\wedge L-dg\wedge KL+g[K,L],\]
 for $K$, $L$ vector-valued one-forms on $J^{1}\pi$ and $g\in C^{\infty}(J^{1}\pi)$.

For the converse, suppose that $R=\alpha\wedge J$, with $\alpha\in T_{v}^{*}$.
Formula (\ref{eq:23'}) implies $\Phi=i_{S} (\alpha\wedge J )= (i_{S}\alpha )J-\alpha\wedge i_{S}J= (i_{S}\alpha )J$.
\end{proof}


\subsection{Lagrangian semisprays}

In this subsection we recall some basic notions about Lagrangian semisprays.

\begin{definition}\null \qquad

1) A smooth function $L\in C^{\infty}(J^{1}\pi)$ is called a \emph{Lagrangian
function}.

2) The Lagrangian $L$ is regular if the $(0,2)$ type tensor with
local components
\begin{gather*}
g_{ij}\big(x^{\alpha},y^{k}\big)=\frac{\partial^{2}L}{\partial y^{i}\partial y^{j}}%\label{eq:tensorg_ij}
\end{gather*}
has rank $n$ on $J^{1}\pi$. The tensor $g=g_{ij}\delta x^{i}\otimes\delta x^{j}$
is called the \emph{metric tensor} of the Lagrangian~$L$.
\end{definition}

\begin{remark}
More exactly, \cite{saunders89}, a function $L\in C^{\infty}(J^{1}\pi)$
is called a Lagrangian density on~$\pi$. If~$\Omega$ is a volume
form on $\mathbb{R}$, the corresponding Lagrangian is the semi-basic
1-form $L\pi_{1}^{*}\Omega$ on~$J^{1}\pi$. Using a f\/ixed volume form
on~$\mathbb{R}$, for example~$dt$, it is natural to consider the
function~$L$ as a~(f\/irst-order) Lagrangian.
\end{remark}

For the particular choice of $dt$ as volume form on $\mathbb{R}$,
the \emph{Poincar\'e--Cartan $1$-form} of the Lagrangian $L$ is
$\theta_{L}:=Ldt+d_{J}L$. The Lagrangian $L$ is regular if and only
if the \emph{Poincar\'e--Cartan $2$-form} $d\theta_{L}$ has maximal
rank $2n$ on $J^{1}\pi$.

For a detailed exposition on the regularity conditions for Lagrangians
see~\cite{Krupkova}.


The \emph{geodesics} of a semispray $S$, given by the system of SODE
(\ref{sode}), coincide with the solutions of the \emph{Euler--Lagrange
equations} \begin{gather*}
\frac{d}{dt}\left(\frac{\partial L}{\partial y^{i}}\right)-\frac{\partial L}{\partial x^{i}}=0%\label{eq:EL}
\end{gather*}
 if and only if \begin{gather}
g_{ij}\left(t,x,\frac{dx}{dt}\right)\left(\frac{d^{2}x^{i}}{dt^{2}}+2G^{i}\left(t,x,\frac{dx}{dt}\right)\right)=\frac{d}{dt}\left(\frac{\partial L}{\partial y^{i}}\right)-\frac{\partial L}{\partial x^{i}}.\label{eq:SEL}\end{gather}


Therefore, for a semispray $S,$ there exists a Lagrangian function
$L$ such that (\ref{eq:SEL}) holds true if and only if $S\left(\frac{\partial L}{\partial y^{i}}\right)-\frac{\partial L}{\partial x^{i}}=0,$
which can be further expressed as
\begin{gather}
\mathcal{L}_{S}\theta_{L}=dL \ \Leftrightarrow \ i_{S}d\theta_{L}=0.\label{eq:LSTL}
\end{gather}


\begin{definition}
A semispray $S$ is called a \emph{Lagrangian semispray} (or a Lagrangian
vector f\/ield) if there exists a Lagrangian function~$L$,
locally def\/ined on~$J^{1}\pi$, that satisf\/ies~(\ref{eq:LSTL}).
\end{definition}

In \cite{B-C-2010} it has been shown that a semispray $S$ is a Lagrangian
semispray if and only if there exists a semi-basic 1-form $\theta\in\Lambda_{v}^{1}$
with $\operatorname{rank}(d\theta)=2n$ on $J^{1}\pi$, such that
$\mathcal{L}_{S}\theta$ is closed. This represents a reformulation,
in terms of semi basic $1$-forms, of the result in terms of $2$-forms
obtained by Crampin et al.\ in~\cite{crampin84a}. The characterization
of Lagrangian higher order semisprays in terms of a closed 2-form
appears also in~\cite{anderso92}.

Based on this result we can obtain the following reformulation in
terms of semi-basic 1-forms of the known Helmholtz conditions \cite[Lemma~4.2, Lemma~4.3,  Theorem~4.5, Theorem~5.1]{B-C-2010}.

\begin{theorem}
A semispray $S$ is a Lagrangian vector field if and only if there
exists a semi-basic $1$-form $\theta\in\Lambda_{v}^{1}$, with $\operatorname{rank}(d\theta)=2n$
on $J^{1}\pi$, such that \begin{gather}
d_{J}\theta=0,\qquad  d_{h}\theta=0.\label{eq:P}
\end{gather}
\end{theorem}
\begin{proof}
In order to make this paper self contained, we give a direct proof
of this theorem.

Suppose that $S$ is a Lagrangian semispray. It results that there
exists a regular Lagrangian~$L$ on~$J^{1}\pi$ with $\mathcal{L}_{S}\theta_{L}=dL$,
or equivalently $i_{S}d\theta_{L}=0$, where $\theta_{L}=Ldt+d_{J}L$
is its Poincar\'e 1-form. Evidently $\theta_{L}$ is a semi-basic 1-form
with $\operatorname{rank}(d\theta_{L})=2n$ on $J^{1}\pi$. We will
prove that $d_{J}\theta_{L}=d_{h}\theta_{L}=0$.

Indeed, $d_{J}\theta_{L}=d_{J}L\wedge dt+Ld_{J}dt+d_{J}^{2}L=i_{J}dL\wedge dt-d_{J\wedge dt}L=i_{J}dL\wedge dt-i_{J\wedge dt}dL=0$.

From the formula $i_{J}\mathcal{L}_{S}-\mathcal{L}_{S}i_{J}=i_{\Gamma-S\otimes dt}$
and $i_{J}\mathcal{L}_{S}d\theta_{L}=0$ we obtain $\mathcal{L}_{S}d_{J}\theta_{L}+i_{\Gamma}d\theta_{L}-i_{S\otimes dt}d\theta_{L}=0$.
We also compute $i_{\Gamma}d\theta_{L}=i_{2h-\operatorname{Id}}d\theta_{L}=2i_{h}d\theta_{L}-2d\theta_{L}=2d_{h}\theta_{L}$.
Therefore $2d_{h}\theta_{L}=i_{S\otimes dt}d\theta_{L}=-i_{S}d\theta_{L}\wedge dt=0$.

Conversely, suppose that there exists a semi-basic $1$-form $\theta\in\Lambda_{v}^{1}$,
with $\operatorname{rank}(d\theta)=2n$ on~$J^{1}\pi$, such that
$d_{J}\theta=0$, $d_{h}\theta=0$. In order to prove that $S$ is
a Lagrangian vector f\/ield, we will f\/irst show that~$\mathcal{L}_{S}\theta=d(i_{S}\theta)$.

The hypothesis $d_{J}\theta=0$ implies $\theta=(i_{S}\theta)dt+d_{J}(i_{S}\theta).$
Indeed, $d_{J}i_{S}+i_{S}d_{J}=\mathcal{L}_{JS}-i_{[S,J]}=i_{h-S\otimes dt-v}\Rightarrow d_{J}(i_{S}\theta)=i_{h}\theta-(i_{S}\theta)dt-i_{v}\theta=\theta-(i_{S}\theta)dt.$

Next, from $d_{h}i_{S}+i_{S}d_{h}=\mathcal{L}_{hS}-i_{[S,h]}$ and
$d_{h}\theta=0$ it results that $d_{h}i_{S}\theta=\mathcal{L}_{S}\theta-i_{\mathbb{F}+J+\Phi}\theta=\mathcal{L}_{S}\theta-i_{\mathbb{F}}\theta$.

{\sloppy From $i_{\mathbb{F}}\theta=i_{\mathbb{F}}\left((i_{S}\theta)dt+d_{J}(i_{S}\theta)\right)=i_{\mathbb{F}}\left(d_{J}(i_{S}\theta)\right)$
and $i_{\mathbb{F}}d_{J}-d_{J}i_{\mathbb{F}}=d_{J\circ\mathbb{F}}-i_{[\mathbb{F},J]}\Rightarrow $ $ i_{\mathbb{F}}\left(d_{J}(i_{S}\theta)\right)=d_{v}(i_{S}\theta)$.
It results that $\mathcal{L}_{S}\theta=d_{h}(i_{S}\theta)+d_{v}(i_{S}\theta)=d(i_{S}\theta)$.

}

Consider $L=i_{S}\theta$. Then $\theta$ is the Poincar\'e--Cartan 1-form
of $L$ and $\mathcal{L}_{S}\theta_{L}=dL$. From $\operatorname{rank}(d\theta)=2n$
on $J^{1}\pi$ it results that $L$ is a regular Lagrangian and $S$
is a Lagrangian vector f\/ield.
\end{proof}

In the next section we discuss the formal integrability of these Helmholtz
conditions using two suf\/f\/icient conditions provided by Cartan--K\"ahler
theorem.

\section{Formal integrability for the nonautonomus inverse problem\\ of the
calculus of variations}\label{section3}

In order to study the integrability conditions of the set of dif\/ferential
equations \eqref{eq:P}, we associate to it a linear  partial dif\/ferential
operator and study its formal integrability, using Spencer's technique.
The approach in this work follows the one developed in \cite{B-M_arxiv}
for studying the projective metrizability problem for autonomous sprays.
For the basic notions of formal integrability theory of linear partial
dif\/ferential operators see \cite{B-M_arxiv,grifone00}.

Consider $T_{v}^{*}$ the vector bundle of semi-basic $1$-forms on
$J^{1}\pi$ and $\Lambda_{v}^{1}$ the module of sections of~$T_{v}^{*}$.
For $\theta\in\Lambda_{v}^{1}$ and $k\geq1$ we denote by $j_{u}^{k}\theta$
the $k$th order jet of $\theta$ at the base point $u$ in $J^{1}\pi$.
The bundle of $k$th order jets of sections of $T_{v}^{*}$ is denoted
by $J^{k}T_{v}^{*}$. The projection $\pi_{0}:J^{k}T_{v}^{*}\rightarrow J^{1}\pi$
is def\/ined by $\pi_{0}(j_{u}^{k}\theta)=u$. If $l>k$, one def\/ines
the projections $\pi_{k}$ as follows: $\pi_{k}(j_{u}^{l}\theta)=j_{u}^{k}\theta$
and $J^{l}T_{v}^{*}$ is also a f\/ibred manifold over $J^{k}T_{v}^{*}$.

If $f_{1},\dots,f_{k}\in C^{\infty}(J^{1}\pi)$ are functions vanishing
at $u\in J^{1}\pi$ and $\theta\in\Lambda_{v}^{1}$, we def\/ine $\epsilon: \mathcal{S}^{k}T^{*}\otimes T_{v}^{*}\overset{}{\longrightarrow}J^{k}T_{v}^{*}$
by $\epsilon(df_{1}\odot\cdots\odot df_{k}\otimes\theta)_{u}=j_{u}^{k}(f_{1}\cdots f_{k}\theta)$,
where $\odot$ is the symmetric product. Then the sequence
\[
0\overset{}{\longrightarrow}\mathcal{S}^{k}T^{*}\otimes T_{v}^{*}\overset{\epsilon}{\longrightarrow}J^{k}T_{v}^{*}\overset{\pi_{k-1}}{\longrightarrow}J^{k-1}T_{v}^{*}\overset{}{\longrightarrow}0\]
 is exact.

Consider the \emph{linear partial differential operator} of order
one\begin{gather}
P: \ \Lambda_{v}^{1}   \rightarrow   \Lambda_{v}^{2}\oplus\Lambda_{v}^{2},\qquad
P   =  \left(d_{J},  d_{h}\right).\label{eq:PDE}
\end{gather}


Remark that $P(\theta)$ can be expressed in terms of f\/irst-order
jets of $\theta$, for any $\theta\in\Lambda_{v}^{1}$, and therefore
it induces a morphism between vector bundles:
\begin{gather*}
p^{0}(P): \  J^{1}T_{v}^{*}   \rightarrow   \Lambda^{2}T_{v}^{*}\oplus\Lambda^{2}T_{v}^{*},\qquad
p^{0}(P)(j_{u}^{1}\theta)   =    P(\theta)_{u},\quad\forall\,\theta\in\Lambda_{v}^{1}.
\end{gather*}


We also consider the \emph{$l$th order jet prolongations} of the
dif\/ferential operator $P$, $l\geq1$, which will be identif\/ied with
the morphisms of vector bundles over $M$,
\begin{gather*}
p^{l}(P): \  J^{l+1}T_{v}^{*}   \rightarrow   J^{l}\left(\Lambda^{2}T_{v}^{*}\oplus\Lambda^{2}T_{v}^{*}\right),\qquad
p^{l}(P)\big(j_{u}^{l+1}\theta\big)   =   j_{u}^{l}\left(P(\theta)\right),\quad \forall\,\theta\in\Lambda_{v}^{1}.
\end{gather*}


Remark that for a semi-basic 1-form $\theta=\theta_{\alpha}\delta x^{\alpha}$,
its f\/irst-order jet $j^{1}\theta=\frac{\delta\theta_{\alpha}}{\delta x^{\beta}}\delta x^{\beta}\otimes\delta x^{\alpha}+\frac{\partial\theta_{\alpha}}{\partial y^{i}}\delta y^{i}\otimes\delta x^{\alpha}$
determines the local coordinates $\left(x^{\alpha},y^{i},\theta_{\alpha},\theta_{\alpha\beta},\theta_{\alpha\underline{i}}\right)$
on $J^{1}T_{v}^{*}$. In this work all contravariant or covariant
indices, related to vertical components of tensor f\/ields will be underlined.

Consider $\theta=\theta_{\alpha}\delta x^{\alpha}$, a semi-basic
$1$-form on $J^{1}\pi$. Then
\begin{gather*}
d\theta   =   \left(\frac{\partial\theta_{i}}{\partial t}-\theta_{j}N_{i}^{j}-\frac{\delta\theta_{0}}{\delta x^{i}}\right)\delta x^{0}\wedge\delta x^{i}+\left(\theta_{i}-\frac{\partial\theta_{0}}{\partial y^{i}}\right)\delta x^{0}\wedge\delta y^{i}\\
\hphantom{d\theta   =}{}
+\frac{1}{2}\left(\frac{\delta\theta_{j}}{\delta x^{i}}-\frac{\delta\theta_{i}}{\delta x^{j}}\right)\delta x^{i}\wedge\delta x^{j}
   +\left(\frac{\partial\theta_{j}}{\partial y^{i}}\right)\delta y^{i}\wedge\delta x^{j},\\
d_{J}\theta   =   \left(\theta_{i}-\frac{\partial\theta_{0}}{\partial y^{i}}\right)\delta x^{0}\wedge\delta y^{i}+\frac{1}{2}\left(\frac{\partial\theta_{i}}{\partial y^{j}}-\frac{\partial\theta_{j}}{\partial y^{i}}\right)\delta x^{j}\wedge\delta x^{i},\\
d_{h}\theta   =   \left(\frac{\partial\theta_{i}}{\partial t}-\theta_{j}N_{i}^{j}-\frac{\delta\theta_{0}}{\delta x^{i}}\right)\delta x^{0}\wedge\delta x^{i}+\frac{1}{2}\left(\frac{\delta\theta_{i}}{\delta x^{j}}-\frac{\delta\theta_{j}}{\delta x^{i}}\right)\delta x^{j}\wedge\delta x^{i}.
\end{gather*}


Using these formulae we obtain
\begin{gather*}
p^{0}(P)\left(j^{1}\theta\right)=
\left(\left(\theta_{i}-\frac{\partial\theta_{0}}{\partial y^{i}}\right)\delta x^{0}\wedge\delta y^{i}+\frac{1}{2}\left(\frac{\partial\theta_{i}}{\partial y^{j}}-\frac{\partial\theta_{j}}{\partial y^{i}}\right)\delta x^{j}\wedge\delta x^{i} , \right.\\
\left.\phantom{p^{0}(P)\left(j^{1}\theta\right)= }{}
\left(\frac{\partial\theta_{i}}{\partial t}-\theta_{j}N_{i}^{j}-\frac{\delta\theta_{0}}{\delta x^{i}}\right)\delta x^{0}\wedge\delta x^{i}+\frac{1}{2}\left(\frac{\delta\theta_{i}}{\delta x^{j}}-\frac{\delta\theta_{j}}{\delta x^{i}}\right)\delta x^{j}\wedge\delta x^{i}\right).
\end{gather*}


The \emph{symbol} of $P$ is the vector bundle morphism $\sigma^{1}(P):  T^{*}\otimes T_{v}^{*}\rightarrow\Lambda^{2}T_{v}^{*}\oplus\Lambda^{2}T_{v}^{*}$
def\/ined by the f\/irst-order terms of $p^{0}(P)$. More exactly, $\sigma^{1}(P)=p^{0}(P)\circ\epsilon$.


For $A\in T^{*}\otimes T_{v}^{*}$, $A=A_{\alpha\beta}\delta x^{\alpha}\otimes\delta x^{\beta}+A_{\underline{i}\beta}\delta y^{i}\otimes\delta x^{\beta}$,
we compute
\begin{gather*}
\sigma^{1}(d_{J})A   =   -A_{\underline{i}0}\delta x^{0}\wedge\delta x^{i}+\frac{1}{2}\big(A_{\underline{j}i}-A_{\underline{i}j}\big)\delta x^{j}\wedge\delta x^{i},%\label{eq:sigma_1_dJ}
\\
\sigma^{1}(d_{h})A   =   \left(A_{0i}-A_{i0}\right)\delta x^{0}\wedge\delta x^{i}+\frac{1}{2}\left(A_{ji}-A_{ij}\right)\delta x^{j}\wedge\delta x^{i}
%\label{eq:sigma_1_dh}
\end{gather*}
 and hence \begin{gather*}
\sigma^{1}(P)A   =   \left(\tau_{J}A , \tau_{h}A\right),\qquad %\label{eq:sigma_1P}\\
\left(\tau_{J}A\right)(X,Y)   =   A(JX,Y)-A(JY,X), \\
\left(\tau_{h}A\right)(X,Y)   =   A(hX,Y)-A(hY,X),
 \end{gather*}
 for $X,Y\in\mathfrak{X}(J^{1}\pi)$. In the above formulae $\tau_{J}$, $\tau_{L}$
are \emph{alternating operators}~\cite{B-M_arxiv}.

\begin{remark}
The alternating operators are def\/ined in general as follows. For $K\in\Psi^{k}(J^{1}\pi)$,
a~vector-valued $k$-form, we consider $\tau_{K}:\Psi^{1}(J^{1}\pi)\otimes\Psi^{l}(J^{1}\pi)\to\Psi^{l+k}(J^{1}\pi)$,
\begin{gather}
  (\tau_{K}B)(X_{1},\dots ,X_{l+k})=
  \frac{1}{l!k!}\sum_{\sigma\in S_{l+k}}\varepsilon(\sigma)B(K(X_{\sigma(1)},\dots ,X_{\sigma(k)}),X_{\sigma(k+1)},\dots ,X_{\sigma(k+l)}),\label{taull}  \end{gather}
 where $X_{1},\dots ,X_{l+k}\in{\mathfrak{X}}(J^{1}\pi)$ and $S_{l+k}$
is the permutation group of $\{1,\dots ,l+k\}$. The restriction of $\tau_{K}$
to $\Psi^{l+1}(J^{1}\pi)$, is a derivation of degree $(k-1)$ and
it coincides with the \emph{inner product} $i_{K}$.
\end{remark}

The \emph{first-order prolongation of the symbol} of $P$ is the vector
bundle morphism
$\sigma^{2}(P):  S^{2}T^{*}\otimes T_{v}^{*}\rightarrow T^{*}\otimes\left(\Lambda^{2}T_{v}^{*}\oplus\Lambda^{2}T_{v}^{*}\right)$
that verif\/ies \[
i_{X}\left(\sigma^{2}(P)B\right)=\sigma^{1}(P)\left(i_{X}B\right),\qquad \forall\, B\in S^{2}T^{*}\otimes T_{v}^{*},\quad \forall\, X\in\mathfrak{X}(J^{1}\pi).
\]


Therefore \begin{gather*}
\sigma^{2}(P)B=\left(\sigma^{2}(d_{J})B , \sigma^{2}(d_{h})B\right),\\
\left(\sigma^{2}(d_{J})B\right)(X,Y,Z)   =   B(X,JY,Z)-B(X,JZ,Y), \\ %\label{eq:sigma_2_dJ}\\
\left(\sigma^{2}(d_{h})B\right)(X,Y,Z)   =   B(X,hY,Z)-B(X,hZ,Y). %\label{eq:sigma_2_dh}
\end{gather*}


In local coordinates we obtain the following formulae.

If $B\in S^{2}T^{*}\otimes T_{v}^{*}$, then it has the local decomposition
\begin{gather}
B   =   B_{\alpha\beta\gamma}\delta x^{\alpha}\otimes\delta x^{\beta}\otimes\delta x^{\gamma}+B_{\underline{i}\alpha\beta}\delta y^{i}\otimes\delta x^{\alpha}\otimes\delta x^{\beta}\nonumber \\
\hphantom{B   =}{} +B_{\alpha\underline{i}\beta}\delta x^{\alpha}\otimes\delta y^{i}\otimes\delta x^{\beta}+B_{\underline{ij}\alpha}\delta y^{i}\otimes\delta y^{j}\otimes\delta x^{\alpha},\label{eq:localB}
 \end{gather}
 with \begin{gather}
B_{\alpha\beta\gamma}=B_{\beta\alpha\gamma},\qquad B_{\underline{i}j\alpha}=B_{i\underline{j}\alpha},\qquad B_{\underline{i}0\alpha}=B_{0\underline{i}\alpha},\qquad B_{\underline{ij}\alpha}=B_{\underline{ji}\alpha}.\label{eq:condB}
\end{gather}


The f\/irst-order prolongation of the symbol of $P$ is given by
\begin{gather*}
\sigma^{2}(d_{J})B   =   B_{\alpha\underline{i}0}\delta x^{\alpha}\otimes\delta x^{i}\wedge\delta x^{0}+B_{\underline{ij}0}\delta y^{i}\otimes\delta x^{j}\wedge\delta x^{0}\nonumber \\
\hphantom{\sigma^{2}(d_{J})B   =}{}
 +\frac{1}{2}\big(B_{\alpha\underline{i}j}-B_{\alpha\underline{j}i}\big)\delta x^{\alpha}\otimes\delta x^{i}\wedge\delta x^{j}+\frac{1}{2}\big(B_{\underline{ij}k}-B_{\underline{ik}j}\big)\delta y^{i}\otimes\delta x^{j}\wedge\delta x^{k},\\ %\label{eq:localsigma2dJ}\\
\sigma^{2}(d_{h})B   =   \frac{1}{2}\left(B_{\alpha\beta\gamma}-B_{\alpha\gamma\beta}\right)\delta x^{\alpha}\otimes\delta x^{\beta}\wedge\delta x^{\gamma}+\frac{1}{2}\left(B_{\underline{i}\alpha\beta}-B_{\underline{i}\beta\alpha}\right)\delta y^{i}\otimes\delta x^{\alpha}\wedge\delta x^{\beta}. \end{gather*}


For each $u\in J^{1}\pi$, we consider
\begin{gather*}
g_{u}^{k}(P)   =   \operatorname{Ker}\sigma_{u}^{k}(P),\qquad k\in\{1,2\},\\
g_{u}^{1}(P)_{e_{1}\dots e_{j}}   =   \{A\in g_{u}^{1}(P)|i_{e_{1}}A=\cdots=i_{e_{j}}A=0\},\qquad j\in\{1,\dots ,n\},
\end{gather*}
 where $\{e_{1},\dots ,e_{n}\}$ is a basis of $T_{u}(J^{1}\pi)$. Such
a basis is called \emph{quasi-regular} if it satisf\/ies
\begin{gather*}
\dim g_{u}^{2}(P)=\dim g_{u}^{1}(P)+\sum_{j=1}^{n}\dim g_{u}^{1}(P)_{e_{1}\dots e_{j}}.%\label{quasir}
\end{gather*}


\begin{definition}
The symbol $\sigma^{1}(\mathcal{P})$ is called \emph{involutive}
at $u$ in $J^{1}\pi$ if there exists a quasi-regular basis of $T_{u}J^{1}\pi$.
\end{definition}

A f\/irst-order jet $j_{u}^{1}\theta\in J^{1}T_{v}^{*}$ is \emph{a~first-order formal solution of~$P$ at~$u$} in $J^{1}\pi$ if $p^{0}(P)(\theta)_{u}=0$.

For $l\geq1$, a $(1+l)$th order jet $j_{u}^{1+l}\theta\in J_{u}^{1+l}T_{v}^{*}$
is \emph{a $(1+l)$th order formal solution of $P$ at $u$} in $J^{1}\pi$
if $p^{l}(P)(\theta)_{u}=0$.

For any $l\geq0$, consider $R_{u}^{1+l}(P)=\ker p_{u}^{l}(P)$ the
space of \emph{$(1+l)$th order formal solutions} of~$P$ at~$u$.
We denote also $\bar{\pi}_{l,u} :  R_{u}^{1+l}(P)\rightarrow R_{u}^{l}(P)$
the restriction of $\pi_{l,u} :  J_{u}^{1+l}(T_{v}^{*})\rightarrow J_{u}^{l}(T_{v}^{*})$
to $R_{u}^{1+l}(P)$.

\begin{definition}
The partial dif\/ferential operator $P$ is called \emph{formally integrable}
at $u$ in $J^{1}\pi$ if $R^{1+l}(P)=\bigcup_{u\in J^{1}\pi}R_{u}^{1+l}(P)$
is a vector bundle over $J^{1}\pi$, for all $l\geq0$, and the map
$\bar{\pi}_{l,u}:R_{u}^{1+l}(P)\rightarrow R_{u}^{l}(P)$ is onto
for all $l\geq1$.
\end{definition}

The f\/ibred submanifold $R^{1}(P)$ of $\pi_{0}:J_{u}^{1}(T_{v}^{*})\rightarrow J^{1}\pi$
is called \emph{the partial differential equation corresponding to
the first-order PDO}~$P$. A solution of the operator $P$ on an open
set $U\subset J^{1}\pi$ is a section $\theta\in\Lambda_{v}^{1}$
def\/ined on $U$ such that $P\theta=0\Leftrightarrow p^{0}(P)(j_{u}^{1}\theta)=0$, $\forall\, u\in U$.

The Cartan--K\"ahler theorem \cite{grifone00} takes the following form
for the particular case of f\/irst-order PDO.

\begin{theorem}
Let $P$ be a first-order linear partial differential operator with
$g^{2}(P)$ a vector bundle over $R^{1}(P)$. If $\overline{\pi}_{1}:R^{2}(P)\to R^{1}(P)$
is onto and the symbol $\sigma^{1}(P)$ is involutive, then $P$ is
formally integrable.
\end{theorem}

\subsection[The involutivity of the symbol of $P$]{The involutivity of the symbol of $\boldsymbol{P}$}

In this subsection we prove that the operator $P$ satisf\/ies one of
the two suf\/f\/icient conditions for formal integrability, provided by
Cartan--K\"ahler theorem: the involutivity of the symbol~$\sigma^{1}(P)$.

\begin{theorem}\label{thm:-involutivity}
The symbol $\sigma^{1}(P)$ of the PDO
$P=(d_{J},d_{h})$ is involutive.
\end{theorem}

\begin{proof}
First we determine $g^{1}(P)=\left\{   A\in T^{*}\otimes T_{v}^{*}\,\vert\,\sigma^{1}(P)A=0\right\} $,
and compute the dimension of its f\/ibers. We obtain
\begin{gather*}
g_{u}^{1}(P)   =   \big\{  A=A_{\alpha\beta}\delta x^{\alpha}\otimes\delta x^{\beta}+A_{\underline{i}\beta}\delta y^{i}\otimes\delta x^{\beta}\,\,\vert\,\, A_{\underline{i}0}=0,\, A_{\underline{j}i}=A_{\underline{i}j},\, A_{\alpha\beta}=A_{\beta\alpha}\big\} .
\end{gather*}
 From $A_{\underline{i}0}=0$ and $A_{\underline{j}i}=A_{\underline{i}j}$
it results that $A_{\underline{i}j}$ contribute with $n(n+1)/2$
components to the dimension of $g_{u}^{1}(P)$, and from $A_{\alpha\beta}=A_{\beta\alpha}$
it follows that $A_{\alpha\beta}$ contribute with $(n+1)(n+2)/2$
components to the dimension of $g_{u}^{1}(P)$. So
\[
\dim g_{u}^{1}(P)=\frac{n(n+1)}{2}+\frac{(n+1)(n+2)}{2}=(n+1)^{2}.
\]


Next we determine $g^{2}(P)=\left\{   B\in S^{2}T^{*}\otimes T_{v}^{*}\,\vert\,\sigma^{2}(P)B=0\right\} $.

If $B\in S^{2}T^{*}\otimes T_{v}^{*}$ has the local components (\ref{eq:localB}),
then $B\in g^{2}(P)$ if and only if the following relations are satisf\/ied:
\begin{gather}
  B_{\alpha\underline{i}0}=0,  \qquad    B_{\underline{ij}0}=0, \qquad
  B_{\alpha\underline{i}j}=B_{\alpha\underline{j}i}, \nonumber \\
  B_{\underline{ij}k}=B_{\underline{ik}j},\qquad
  B_{\alpha\beta\gamma}=B_{\alpha\gamma\beta}, \qquad B_{\underline{i}\alpha\beta}=B_{\underline{i}\beta\alpha}.\label{eq:eq:cond_B_in_g2}
  \end{gather}


From the relations \eqref{eq:condB} and \eqref{eq:eq:cond_B_in_g2}
it results that $B\in g^{2}(P)$ if and only if its local components
$  B_{\alpha\beta\gamma}$, $B_{\underline{i}jk}$, $B_{i\underline{j}k}$, $B_{\underline{ij}k}$
are totally symmetric and the rest are vanishing. Therefore $B_{\alpha\beta\gamma}$
contribute with $(n+1)(n+2)(n+3)/6$ components to the dimension of
$g_{u}^{2}(P)$, and $B_{\underline{i}jk}$,  $B_{\underline{ij}k}$
with $n(n+1)(n+2)/6$ components each of them. It results
\[
\dim g_{u}^{2}(P)=\frac{(n+1)(n+2)(n+3)}{6}+2\frac{n(n+1)(n+2)}{6}=\frac{(n+1)^{2}(n+2)}{2}.
\]


Consider \[
\mathcal{B}=\left\{ h_{0}=S,\, h_{1},\,\dots,h_{n},\, v_{1}=Jh_{1},\,\dots,\, v_{n}=Jh_{n}\right\} \]
 a basis in $T_{u}J^{1}\pi$ with $h_{0}=S$, $h_{1},\dots,h_{n}$
horizontal vector f\/ields. For any $A\in g^{1}(P)$, we denote \[
A(h_{\alpha},h_{\beta})=a_{\alpha\beta},\qquad A(v_{i},h_{\alpha})=b_{\underline{i}\alpha}.\]


Because $A$ is semi-basic in the second argument it follows that
these are the only components of $A$.

Since $A\in\operatorname{Ker}\sigma^{1}(d_{J})$ it follows that $A_{\underline{i}0}=0$, $A_{\underline{i}j}=A_{\underline{j}i}$ $\Rightarrow$ $b_{\underline{i}0}=0$, $b_{\underline{i}j}=b_{\underline{j}i}$.
Because $A\in\operatorname{Ker}\sigma^{1}(d_{h})$ it results that
$A_{0i}=A_{i0}$, $A_{ij}=A_{ji}$ and hence $a_{0i}=a_{i0}$, $a_{ij}=a_{ji}$.

Consider $j\in \{1,\dots, n\} $ arbitrarily f\/ixed and
\begin{gather*}
\mathcal{\tilde{B}}=\big\{ e_{0}=S+h_{j}+v_{n},\, e_{1}=h_{1},\, e_{2}=h_{2}+v_{1}, \, \dots,\\
 \hphantom{\mathcal{\tilde{B}}=\big\{}{} e_{i}=h_{i}+v_{i-1},\, \dots,\, e_{n}=h_{n}+v_{n-1},\, v_{1},\, \dots,\, v_{n}\big\}
\end{gather*}
 a new basis in $T_{u}J^{1}\pi$. If we denote \[
A(e_{\alpha},e_{\beta})=\tilde{a}_{\alpha\beta},\qquad A(v_{i},e_{\alpha})=\tilde{b}_{\underline{i}\alpha},\]
 a simple computation and the fact that $A$ is semi-basic in the
second argument determine \begin{gather*}
\tilde{a}_{00}   =   a_{00}+2a_{0j}+a_{jj}+b_{\underline{n}j},\\
\tilde{a}_{ik}   =   a_{ik}+b_{\underline{i-1},k}\neq\tilde{a}_{ki}=a_{ki}+b_{\underline{k-1},i},\\
\tilde{a}_{i0}   =   a_{i0}+a_{ij}+b_{\underline{i-1},j}\neq\tilde{a}_{0i}=a_{0i}+a_{ji}+b_{\underline{n}i},\\
\tilde{b}_{\underline{i}k}   =b_{\underline{i}k}   =\tilde{b}_{\underline{k}i},\qquad
\tilde{b}_{\underline{i}0}   =   b_{\underline{i}j}.
\end{gather*}


It can be seen that all the independent components of $A$ in the
basis $\mathcal{B}$ can be obtained from the components of $A$ in
the basis $\mathcal{\tilde{B}}$, and hence we can use the later for
determining the dimensions of $(g_{u}^{1})_{e_{0} e_{1}\dots e_{k}}$.

If $A\in(g_{u}^{1})_{e_{0}}$ it results $\tilde{a}_{0\alpha}=0$,
so using this new basis we impose $n+1$ supplementary independent
restrictions. It follows that $\dim(g_{u}^{1})_{e_{0}}=(n+1)^{2}-(n+1)=(n+1)n$.

If $A\in(g_{u}^{1})_{e_{0} e_{1}}$ it results that together with
the previous restrictions we impose also $\tilde{a}_{1\alpha}=0$,
so another independent $n+1$ restrictions. Hence $\dim (g_{u}^{1})_{e_{0} e_{1}}=(n+1)n-(n+1)=(n+1)(n-1)$.

In general $\dim (g_{u}^{1})_{e_{0} e_{1} \dots e_{k}}=(n+1)(n-k)$, $\forall\, k\in \{1,\dots,n\}$.
Hence $\dim (g_{u}^{1})_{e_{0} e_{1} \dots e_{n}}=0\Rightarrow\dim (g_{u}^{1})_{e_{0} \dots e_{n} v_{1} \dots v_{k}}=0$, $\forall\, k\in \{1,\dots,n\}$,
\begin{gather*}
\dim (g_{u}^{1})+\sum_{k=0}^{n}\dim (g_{u}^{1})_{e_{0} e_{1} \dots e_{k}}
   =   (n+1)^{2}+(n+1)[n+n-1+\cdots+1]\\
\hphantom{\dim (g_{u}^{1})+\sum_{k=0}^{n}\dim (g_{u}^{1})_{e_{0} e_{1} \dots e_{k}}}{}   =   \frac{(n+1)^{2}(n+2)}{2}=\dim g_{u}^{2}(P).
   \end{gather*}


We proved that $\tilde{B}$ is a a quasi-regular basis, hence the
symbol $\sigma^{1}(P)$ is involutive.
\end{proof}




\subsection{First obstruction to the inverse problem}

In this subsection we determine necessary and suf\/f\/icient conditions
for $\bar{\pi}_{1}$ to be onto. We will obtain only one obstruction
for the integrability of the operator $P$. The obstruction is due
to the curvature tensor of the nonlinear connection induced by the
semispray.

\begin{theorem}
\label{thm:A-first-order}A first-order formal solution $\theta\in\Lambda_{v}^{1}$
of the system $d_{J}\theta=0$, $d_{h}\theta=0$ can be lifted into
a second-order solution, which means that $\bar{\pi}_{1} :  R^{2}(P)\rightarrow R^{1}(P)$
is onto, if and only if \[
d_{R}\theta=0,\]
 where $R$ is the curvature tensor \eqref{eq:19}.
\end{theorem}
\begin{proof}
We use a known result from \cite[Proposition~1.1]{grifone00}.

If $\mathcal{K}$ is the cokernel of $\sigma^{2}(P)$,
\begin{gather*}
\mathcal{K}=\frac{T^{*}\otimes\left(\Lambda^{2}T_{v}^{*}\oplus\Lambda^{2}T_{v}^{*}\right)}{{\normalcolor \operatorname{Im}}\sigma^{2}(P)},%\label{eq:K}
\end{gather*}
 there exists a morphism $\varphi :R^{1}(P)\rightarrow\mathcal{K}$
such that the sequence \[
R^{2}(P)\overset{\bar{\pi}_{1}}{\longrightarrow}R^{1}(P)\overset{\varphi}{\longrightarrow}\mathcal{K}\]
 is exact. In particular $\bar{\pi}_{1}$ is onto if and only if $\varphi=0$.

After def\/ining $\varphi$, we will prove that for $\theta\in\Lambda_{v}^{1}$,
with $j_{u}^{1}\theta\in R_{u}^{1}(P)$, a f\/irst-order formal solution
of $P$ at $u\in J^{1}\pi$, we have that $\varphi_{u}\theta=0$ if
and only if $\left(d_{R}\theta\right)_{u}=0$.

The construction of the morphism $\varphi$ is represented in the
next diagram by dashed arrows. We denote $F=\Lambda^{2}T_{v}^{*}\oplus\Lambda^{2}T_{v}^{*}$.
%$$  {\diagram & 0 \dto & 0 \dto & 0 \dto \\ 0\rto & g^2(P) \rto \dto & S^{2}T^*\otimes T^*_v \rto^{\sigma^2(P)} \dto^{\varepsilon} & T^* \otimes F %\rto \rdashed<1ex>|>{\tip}^{\tau} \dto^{\varepsilon} &  \it{K} \rto & 0 \\ 0 \rto & R^{2}(P) \rto^{i} \dto^{\overline {\pi}_{1}} & J^{2} T^*_v \rto %\dto^{\pi_{1}} \rdashed<1ex>|>{\tip}^{p^1(P)} & J^1 F \dto^{\pi_0} \udashed<1ex>|>{\tip}^{p^0(\nabla)} \\ 0 \rto & R^1(P) \rto\rdashed<1ex>|>{\tip}^i %& J^1T^*_v \rto^{p^0(P)} \udashed<1ex>|>{\tip} \dto & F \dto \\ & & 0 & 0 \enddiagram}
%$$
$$
\includegraphics{Diagram}
$$



Remark that $\dim T^{*}=2n+1$, $\dim T_{v}^{*}=n+1$, $\dim\Lambda^{2}T_{v}^{*}=\frac{(n+1)n}{2}$, $\dim\mathcal{S}^{2}T^{*}=\frac{(2n+1)(2n+2)}{2}$,
$\dim T^{*}\otimes\left(\Lambda^{2}T_{v}^{*}\oplus\Lambda^{2}T_{v}^{*}\right)=(2n+1)n(n+1)$.
Therefore
\begin{gather*}
\dim\mathcal{K}   =   \dim\left[T^{*}\otimes\left(\Lambda^{2}T_{v}^{*}\oplus\Lambda^{2}T_{v}^{*}\right)\right]-\dim\left(\operatorname{Im}\sigma^{2}(P)\right)\\
\hphantom{\dim\mathcal{K}}{} =   \dim\left[T^{*}\otimes\left(\Lambda^{2}T_{v}^{*}\oplus\Lambda^{2}T_{v}^{*}\right)\right]-\left[\dim\left(\mathcal{S}^{2}T^{*}\otimes T_{v}^{*}\right)-\dim\left(\ker\sigma^{2}(P)\right)\right]\\
\hphantom{\dim\mathcal{K}}{}
 =  \frac{(n-1)n(n+1)}{2}=3 \begin{pmatrix}
n+1\\
3\end{pmatrix}.\end{gather*}


It results from this that \[
\mathcal{K}\simeq\oplus^{(3)}\Lambda^{3}T_{v}^{*}.\]


Next we def\/ine $\tau:  T^{*}\otimes\left(\Lambda^{2}T_{v}^{*}\oplus\Lambda^{2}T_{v}^{*}\right)\rightarrow\oplus^{(3)}\Lambda^{3}T_{v}^{*}$
such as the next sequence is exact:
 \begin{gather}
0\rightarrow g^{2}(P)\overset{i}{\longrightarrow}\mathcal{S}^{2}T^{*}\otimes T_{v}^{*}\overset{\sigma^{2}(P)}{\longrightarrow}T^{*}
\otimes\left(\Lambda^{2}T_{v}^{*}\oplus\Lambda^{2}T_{v}^{*}\right)\overset{\tau}{\longrightarrow}\oplus^{(3)}\Lambda^{3}T_{v}^{*}\rightarrow0.
\label{eq:exact_tau}\end{gather}


For $B_{1},  B_{2}\in T^{*}\otimes\Lambda^{2}T_{v}^{*}$, we def\/ine
$\tau(B_{1},B_{2})=\left(\tau_{1}(B_{1},B_{2}) , \tau_{2}(B_{1},B_{2}) , \tau_{3}(B_{1},B_{2})\right)$,
where $\tau_{i}:  T^{*}\otimes\left(\Lambda^{2}T_{v}^{*}\oplus\Lambda^{2}T_{v}^{*}\right)\rightarrow\Lambda^{3}T_{v}^{*}$, $ i\in \{1,2,3\}$,
are given by
\begin{gather*}
\tau_{1}(B_{1},B_{2})=\tau_{J}B_{1},\qquad
\tau_{2}(B_{1},B_{2})=\tau_{h}B_{2},\qquad \tau_{3}(B_{1},B_{2})=\tau_{h}B_{1}+\tau_{J}B_{2}.%\label{eq:tau_123}
\end{gather*}


Using the def\/inition (\ref{taull}) of the alternating operators $\tau_{J}$, $\tau_{h}$,
we prove that $\tau\circ\sigma^{2}(P)=0$.

Indeed, using that any $B\in\mathcal{S}^{2}T^{*}\otimes T_{v}^{*}$
is symmetric in the f\/irst two arguments, it follows that $\left(\tau\circ\sigma^{2}(P)\right)(B)=\tau\left(\sigma^{2}(d_{J})B, \sigma^{2}(d_{h})B\right)=\left(\tau_{J}\sigma^{2}(d_{J})B, \tau_{h}\sigma^{2}(d_{h})B, \tau_{h}\sigma^{2}(d_{J})B+\tau_{J}\sigma^{2}(d_{h})B\right)=0$, $\forall\, B\in\mathcal{S}^{2}T^{*}\otimes T_{v}^{*}$.
For example,\begin{gather*}
\tau_{J}\left(\sigma^{2}(d_{J})B\right)(X,Y,Z)  \\
 \qquad{} =  \left[\sigma^{2}(d_{J})B\left(JX,Y,Z\right)-\sigma^{2}(d_{J})B\left(JY,X,Z\right)+\sigma^{2}(d_{J})B\left(JZ,X,Y\right)\right]\\
  \qquad{}{} =   [B(JX,JY,Z)-B(JX,JZ,Y)-B(JY,JX,Z)+B(JY,JZ,X)\\
   \qquad\quad{}{}   +B(JZ,JX,Y)-B(JZ,JY,X)]=0,\qquad \forall\, X,Y,Z\in\mathfrak{X}(J^{1}\pi).
 \end{gather*}


The relation $\tau\circ\sigma^{2}(P)=0$ implies that $\operatorname{Im}(\sigma^{2}(P))\subseteq\operatorname{Ker}\tau$.
Using that $\tau$ is onto ($\tau_{J}$,~$\tau_{h}$~are both onto)
it results that $\dim\left[\operatorname{Im}(\sigma^{2}(P))\right]=\dim\left(\operatorname{Ker}\tau\right)$
and hence $\operatorname{Im}(\sigma^{2}(P))=\operatorname{Ker}\tau$
and the sequence~(\ref{eq:exact_tau}) is exact.

The last step before def\/ining $\varphi :R^{1}(P)\rightarrow\mathcal{K}$
is to consider a linear connection $\nabla$ on $J^{1}\pi$ such that
$\nabla J=0$. It means that $\nabla$ preserve semi-basic forms and
$\nabla$ can be considered as a~connection in the f\/iber bundle $\Lambda^{2}T_{v}^{*}\oplus\Lambda^{2}T_{v}^{*}\rightarrow J^{1}\pi$.
As a f\/irst-order PDO we can identify $\nabla$ with the bundle morphism
$p^{0}(\nabla):  J^{1}\left(\Lambda^{2}T_{v}^{*}\oplus\Lambda^{2}T_{v}^{*}\right)\rightarrow T^{*}\otimes\left(\Lambda^{2}T_{v}^{*}\oplus\Lambda^{2}T_{v}^{*}\right)$.

We will also use two derivations of degree~1 introduced in~\cite{B-M_arxiv},
def\/ined by $\mathcal{D}_{J}=\tau_{J}\nabla$, $\mathcal{D}_{h}=\tau_{h}\nabla.$
Both derivations $\mathcal{D}_{J}$, $\mathcal{D}_{h}$ preserve semi-basic
forms and $d_{J}-\mathcal{D}_{J}$, $ d_{h}-\mathcal{D}_{h}$ are algebraic
derivations. It means that if $\omega\in\Lambda^{k}\left(J^{1}\pi\right)$
vanishes at some point $u\in J^{1}\pi$, then $\left(\mathcal{D}_{J}\omega\right)_{u}=\left(d_{J}\omega\right)_{u}$
and $\left(\mathcal{D}_{h}\omega\right)_{u}=\left(d_{h}\omega\right)_{u}$
\cite[Lemma~2.1]{B-M_arxiv}.

Now we are able to def\/ine $\varphi :R^{1}(P)\rightarrow\mathcal{K}$
such that the sequence \[
R^{2}(P)\overset{\bar{\pi}_{1}}{\longrightarrow}R^{1}(P)\overset{\varphi}{\longrightarrow}\mathcal{K}\]
 is exact.

Let $\theta\in\Lambda_{v}^{1}$ such that $j_{u}^{1}\theta\in R_{u}^{1}(P)\subset J_{u}^{1}T_{v}^{*}$
is a f\/irst-order formal solution of $P$ at $u\in J^{1}\pi$, which
means that $\left(d_{J}\theta\right)_{u}=\left(d_{h}\theta\right)_{u}=0$.

Consider \[
\varphi_{u}\theta=\tau_{u}\nabla P\theta=\tau_{u}\left(\nabla d_{J}\theta, \nabla d_{h}\theta\right).\]


Using the fact that $d_{J}-\tau_{J}\nabla$ and $d_{h}-\tau_{h}\nabla$
are algebraic derivations and $\left(d_{J}\theta\right)_{u}=\left(d_{h}\theta\right)_{u}=0$
it results $d_{J}\left(d_{J}\theta\right)_{u}=\tau_{J}\nabla\left(d_{J}\theta\right)_{u}$
and $d_{h}\left(d_{h}\theta\right)_{u}=\tau_{h}\nabla\left(d_{h}\theta\right)_{u}$.

We will compute the three components of the map $\varphi$. It follows
that \begin{gather*}
\tau_{1}\left(\nabla P\theta\right)_{u}   =   \tau_{1}\left(\nabla d_{J}\theta, \nabla d_{h}\theta\right)_{u}=\tau_{J}\left(\nabla d_{J}\theta\right)_{u}  \\
\hphantom{\tau_{1}\left(\nabla P\theta\right)_{u}}{}
 =\left(d_{J}^{2}\theta\right)_{u}=\frac{1}{2}\left(d_{[J,J]}\theta\right)_{u}=-\left(d_{J\wedge dt}\theta\right)_{u}=-\left(d_{J}\theta\right)_{u}\wedge dt=0, %\label{eq:tau_1nablaPtheta}
 \\
\tau_{2}\left(\nabla P\theta\right)_{u}   =   \tau_{2}\left(\nabla d_{J}\theta, \nabla d_{h}\theta\right)_{u}=\tau_{h}\left(\nabla d_{h}\theta\right)_{u}
   =   \left(d_{h}^{2}\theta\right)_{u}=\frac{1}{2}\left(d_{[h,h]}\theta\right)_{u}=\left(d_{R}\theta\right)_{u},%\label{eq:obstructiond_R}
   \end{gather*}
 where $R$ is given by (\ref{eq:19}),
\begin{gather*}
\tau_{3}\left(\nabla P\theta\right)_{u}   =   \tau_{3}\left(\nabla d_{J}\theta, \nabla d_{h}\theta\right)_{u}=\tau_{h}\left(\nabla d_{J}\theta\right)_{u}+\tau_{J}\left(\nabla d_{h}\theta\right)_{u}\nonumber \\
\hphantom{\tau_{3}\left(\nabla P\theta\right)_{u}}{}
   =   \left(d_{h}d_{J}\theta\right)_{u}+\left(d_{J}d_{h}\theta\right)_{u}=\left(d_{[h,J]}\theta\right)_{u}=0.%\label{eq:tau_3nablaPtheta}
 \end{gather*}
Hence $\varphi=0$ if and only if $d_{R}\theta=0$.
\end{proof}
